% !TEX root = ./ECEN-629 - Homework 3 - Brody Jordan.tex

\documentclass[11pt]{article}

\usepackage[
    assignmentDueDate={October 21, 2025 at 12:45 pm}
]{C:/Users/bljor/Sync/Courses/assignment}

\usepackage{algorithm}
\usepackage[noend]{algpseudocode}
\DeclareMathOperator{\minimize}{minimize:}
\DeclareMathOperator{\subjectto}{subject\ to:}
\usetikzlibrary{arrows,intersections,math,patterns,fillbetween}
\usepgfplotslibrary{fillbetween}

\begin{document}
\makeAssignmentTitle

% ----- PROBLEM 1 ------ %
\begin{homeworkProblem}
    Suppose $f : \R[n] \to \R$ is convex, and define
    \[
        g(x, t) = f(x/t), \quad \dom g = \{(x, t) \mid x/t \in \dom f, \ t > 0\}.
    \]
    Show that $g$ is quasiconvex, and then write down the optimization procedure similarly to that on page 17
    of the slides of Chapter 4, specialized to this problem. \\

    Hint: consider the function $t \cdot f(x/t)$. \\

    \solution

    From Chapter 3, we know that $f : \R[n] \to \R$ is quasiconvex if $\dom f$ is convex and the sublevel sets
    \[
        S_\alpha = \{x \in \dom f \mid f(x) \leq \alpha \}
    \] are convex for all $\alpha$. For $g$:
    \[
        S_\alpha = \{(x, t) \in \dom g \mid g(x, t) \leq \alpha \} = \{(x, t) \mid f(x/t) \leq \alpha, \ t > 0\}.
    \]
    Since $t > 0$, the sublevel sets of $g$ are scaled sublevel sets of $f$, which are convex.
    Thus, $g$ is quasiconvex. \\

    \begin{algorithm}[H]
        \caption{Quasiconvex Optimization}
        \begin{algorithmic}[1]
            \Require $l \le p^\star$, $u \ge p^\star$, tolerance $\varepsilon > 0$
            \Repeat
            \State $t \gets (l + u)/2$
            \State Solve the convex feasibility problem:
            \[
                f(x/t) \leq \alpha, \quad t > 0
            \]
            \If{feasible}
            \State $u \gets t$
            \Else
            \State $l \gets t$
            \EndIf
            \Until{$u - l \le \varepsilon$}
            \State \Return $t^\star \approx u$
        \end{algorithmic}
    \end{algorithm}


\end{homeworkProblem}

\pagebreak

% ----- PROBLEM 2 ------ %
\begin{homeworkProblem}
    Some non-convex problems can also be easy to solve. Find the optimal solution for the
    following problems in closed form.
    \begin{enumerate}[label=(\alph*), topsep=5pt, itemsep=5pt]
        \item For $x \in \R[n]$
              \begin{align*}
                  \minimize  & \quad  \norm{x}_2                 \\
                  \subjectto & \quad  10 \geq \norm{x}_1 \geq 1.
              \end{align*}
        \item Let $x = (x_1, x_2) \in \R[2]$
              \begin{align*}
                  \minimize  & \quad  x_1 + x_2                          \\
                  \subjectto & \quad  x_1 + 3x_2 = 5;                    \\
                             & \quad  (x_1 - 3)^2 + (x_2 - 4)^2 \geq 25. \\
              \end{align*}
    \end{enumerate}

    \solution

    \textbf{Part (a):} The objective function $f_0(x) = \sqrt{\sum_{i=1}^{n} |x_i|^2}$.
    Choosing the lower bound of the inequality constraint requires $\norm{x}_1 = 1$.
    The $\ell_2$ norm is minimal when each component of $x_i \leq 1$. Thus we aim for
    each component to be minimal, but sum to 1. Thus, $x_i^* = 1/n \text{ for } i = 1, 2, \dots, n$.\\

    \textbf{Part (b):} This problem is unbounded below, so the optimal solution is $-\infty$.

\end{homeworkProblem}

\pagebreak

% ----- PROBLEM 3 ------ %
\begin{homeworkProblem}
    Huber's penalty function is defined as
    \[
        \phi(u) =
        \begin{cases}
            u^2,         & |u| \leq M, \\
            M(2|u| - M), & |u| > M.
        \end{cases}
    \]
    \begin{enumerate}[label=(\alph*), topsep=5pt, itemsep=5pt]
        \item Is the function $\phi(u)$ convex in general?
        \item Show that the minimization problem
              \[
                  \minimize \quad \sum_{i=1}^{n} \phi(a_i^T x - b_i)
              \]
              is equivalent to the following problem
              \begin{align*}
                  \minimize  & \quad  \sum_{i=1}^{n} (u_i^2 + 2M v_i)     \\
                  \subjectto & \quad  -u - v \preceq Ax - b \preceq u + v \\
                             & \quad  0 \preceq u \preceq M\mathbf{1}     \\
                             & \quad  v \succeq 0.
              \end{align*}
        \item A general form of Huber's penalty function with parameters $(\alpha, \beta, \gamma)$ is
              \[
                  \phi(u) = \begin{cases}
                      \alpha \cdot u^2,         & |u| \leq M, \\
                      \beta M (\gamma |u| - M), & |u| > M .
                  \end{cases}
              \]
              For any fixed value of $\alpha > 0$, find the valid assignment of $(\beta, \gamma)$ in terms of $\alpha$
              such that the penalty function is convex. \\
    \end{enumerate}

    \solution

    \textbf{Part (a):} Yes, $\phi(u)$ is convex. Quadratic functions are convex, and $M(2|u| - M)$ is also convex. \\

    \textbf{Part (b):} Here, expanding the objective function gives
    \[
        \phi(a_i^T x - b_i) = \begin{cases}
            (a_i^T x - b_i)^2,       & |a_i^T x - b_i| \leq M, \\
            M(2|a_i^T x - b_i| - M), & |a_i^T x - b_i| > M.
        \end{cases}
    \] where the constraints become
    \[
        -u - v \preceq Ax - b \preceq u + v, \quad 0 \preceq u \preceq M\mathbf{1}, \quad v \succeq 0.
    \]
    \textbf{Part (c):} By continuity, we require \[
        \alpha \cdot M^2 = \beta M (\gamma M - M) \implies \beta = \frac{\alpha M^2}{M(\gamma M - M)} = \frac{\alpha}{\gamma - 1}.
    \]

\end{homeworkProblem}

\pagebreak

% ----- PROBLEM 4 ------ %
\begin{homeworkProblem}
    Write the following problems as equivalent LPs:
    \begin{enumerate}[label=(\alph*), topsep=5pt, itemsep=5pt]
        \item \begin{align*}
                  \minimize  & \quad \norm{x}_\infty        \\
                  \subjectto & \quad \norm{Ax - b}_1 \leq 1
              \end{align*} where $A \in \R[m \times n]$ and $b \in \R[m]$.
        \item \begin{align*}
                  \minimize  & \quad \norm{x}_\infty             \\
                  \subjectto & \quad \norm{Ax - b}_\infty \leq 1
              \end{align*} where $A \in \R[m \times n]$ and $b \in \R[m]$.
        \item \begin{align*}
                  \minimize  & \quad \norm{x}_1                  \\
                  \subjectto & \quad \norm{Ax - b}_\infty \leq 1
              \end{align*} where $A \in \R[m \times n]$ and $b \in \R[m]$.
        \item \begin{align*}
                  \minimize  & \quad \norm{x}_1             \\
                  \subjectto & \quad \norm{Ax - b}_1 \leq 1
              \end{align*} where $A \in \R[m \times n]$ and $b \in \R[m]$.
    \end{enumerate}
    For the problems above, what if we exchange the objective function and the constraint, can you still write down the problems as LPs? \\

    \solution

    \textbf{Part (a):}
    \begin{align*}
        \minimize  & \quad t                                                  \\
        \subjectto & \quad \mathbf{1} \cdot -t  \preceq x \preceq \mathbf{1}t \\
                   & \quad -\mathbf{1}a \preceq Ax - b \preceq \mathbf{1}a    \\
                   & \quad \mathbf{1}^T a \preceq 1, \, \quad a \succeq 0.
    \end{align*}

    \pagebreak

    \textbf{Part (b):}
    \begin{align*}
        \minimize  & \quad t                                                  \\
        \subjectto & \quad \mathbf{1} \cdot -t  \preceq x \preceq \mathbf{1}t \\
                   & \quad -\mathbf{1} \preceq Ax - b \preceq \mathbf{1}.     \\
    \end{align*}

    \textbf{Part (c):}
    \begin{align*}
        \minimize  & \quad \mathbf{1} t                                   \\
        \subjectto & \quad -t \preceq x \preceq t                         \\
                   & \quad -\mathbf{1} \preceq Ax - b \preceq \mathbf{1}. \\
    \end{align*}

    \textbf{Part (d):}
    \begin{align*}
        \minimize  & \quad \mathbf{1} t                                      \\
        \subjectto & \quad -t \preceq x \preceq t                            \\
                   & \quad -\mathbf{1} a \preceq Ax - b \preceq \mathbf{1} a \\
                   & \quad \mathbf{1}^T a \preceq 1, \, \quad a \succeq 0.   \\
    \end{align*}



\end{homeworkProblem}

\pagebreak

% ----- PROBLEM 5 ------ %
\begin{homeworkProblem}
    Consider the optimization problem in $\R[2]$
    \begin{align*}
        \minimize  & \quad  f_0(x_1, x_2)                    \\
        \subjectto & \quad  (x_1 - 1)^2 + (x_2 - 2)^2 \leq 2 \\
                   & \quad  2x_1 + x_2 \geq 3                \\
                   & \quad  x_1 \geq 0, \ x_2 \geq 0.
    \end{align*}
    First make a sketch of the feasible set. For each of the following objective functions, give
    the optimal set and the optimal value, using two methods: geometrically and the first-order
    optimality condition. Note that you should not use KKT directly here, but the first-order
    condition in Chapter 4.
    \begin{enumerate}[label=(\alph*), topsep=5pt, itemsep=5pt]
        \item $f_0(x_1, x_2) = x_1 + x_2$;
        \item $f_0(x_1, x_2) = x_1^2 + x_2^2$. \newline
    \end{enumerate}

    \solution
    \begin{figure}[hbpt!]
        \centering

        \begin{tikzpicture}[scale=1.5]
            % Axes
            \draw[->, thick] (0,0) -- (4,0) node[right] {$x_1$};
            \draw[->, thick] (0,0) -- (0,4) node[above] {$x_2$};
            % draw axis ticks
            \foreach \x in {1,2,3}
            \draw (\x,0.1) -- (\x,-0.1) node[below] {\x};
            \foreach \y in {1,2,3}
            \draw (0.1,\y) -- (-0.1,\y) node[left] {\y};
            \draw(3/2, 0.1) -- (3/2, -0.1) node[below] {$\frac{3}{2}$};


            \begin{scope}
                % Clip to the circle centered at (1,2) with radius sqrt(2) ≈ 1.414
                \clip (1,2) circle (1.414);
                % Clip to the region above the line 2x_1 + x_2 = 3
                % To represent 2x_1 + x_2 > 3, use a large rectangle above the line
                \clip (0,3) -- (3/2,0) -- (4,0) -- (4,4) -- (0,4) -- cycle;
                % Fill the clipped region
                \fill[color=green!60, opacity=0.6] (0,0) rectangle (4,4);
            \end{scope}
            % Lines
            \draw[thick,blue]  (0,3) -- (3/2,0) node[right] {};

            % Vertices (extreme points)
            \filldraw[color=black!60, opacity=0.3, fill=red!5, very thick,](1,2) circle (1.414);
            \filldraw[black] (1.2,0.6) circle (1pt) node[right] {$(1.2, 0.6)$};

        \end{tikzpicture}

    \end{figure}

    \textbf{Part (a):} Geometrically, the optimal solution lies at the lower intersection, which is $x^* = (1.2, 0.6)$, $f(x^*) = 1.2 + 0.6 = 1.8$. \\

    \textbf{Part (b):} Geometrically, the optimal solution lies at the point closest to the origin, with preference for values less than $1$.
    In this case that point is also $x^* = (1.2, 0.6)$, $f(x^*) = 1.2^2 + 0.6^2 = 1.8$.
\end{homeworkProblem}

\pagebreak

% ----- PROBLEM 6 ------ %
\begin{homeworkProblem}
    Consider the optimization problem
    \begin{align*}
        \minimize  & \quad  f_0(x_1, x_2)             \\
        \subjectto & \quad  x_1 + x_2 \geq 2          \\
                   & \quad  2x_1 + x_2 \geq 3         \\
                   & \quad  x_1 \geq 0, \ x_2 \geq 0.
    \end{align*}
    First make a sketch of the feasible set. For each of the following objective functions, give
    the optimal set and the optimal value, using two methods: geometrically and the first-order
    optimality condition.
    \begin{enumerate}[label=(\alph*), topsep=5pt, itemsep=5pt]
        \item $f_0(x_1, x_2) = 3x_1 + 2x_2$;
        \item $f_0(x_1, x_2) = \max\{x_1, x_2\}$. \newline
    \end{enumerate}

    \solution

    \begin{figure}[hbpt!]
        \centering
        \begin{tikzpicture}[scale=1.5]
            % Axes
            \draw[->, thick] (0,0) -- (4,0) node[right] {$x_1$};
            \draw[->, thick] (0,0) -- (0,4) node[above] {$x_2$};

            % Lines
            \draw[thick,blue]  (0,2) -- (2,0) node[right] {};
            \draw[thick,red]   (0,3) -- (1.5,0) node[right] {};

            % Feasible region shading
            \fill[green!20]
            (0,3) -- (1,1) -- (2,0) -- (3.5,0) -- (3.5,3.5) -- (0,3.5) -- cycle;

            % Feasible region boundary (thicker)
            \draw[thick,green!60!black] (0,3) -- (1,1) -- (2,0);

            % Vertices (extreme points)
            \filldraw[black] (0,3) circle (1pt) node[above left] {$(0,3)$};
            \filldraw[black] (1,1) circle (1pt) node[above right] {$(1,1)$};
            \filldraw[black] (2,0) circle (1pt) node[below right] {$(2,0)$};
        \end{tikzpicture}

    \end{figure}

    \textbf{Part (a):} Geometrically, it is clear that the optimal solution is $x = (1, 1)$, $f(x) = 3(1) + 2(1) = 5$. \\

    \textbf{Part (b):} Geometrically, it is clear that the optimal solution is also $x = (1, 1)$, $f(x) = \max\{1, 1\} = 1$.

\end{homeworkProblem}

\pagebreak

% ----- PROBLEM 7 ------ %
\begin{homeworkProblem}
    Give explicit solutions for the following LP problems over $\R[n]$
    \begin{enumerate}[label=(\alph*), topsep=5pt, itemsep=5pt]
        \item \begin{align*}
                  \minimize  & \quad  c^T x         \\
                  \subjectto & \quad  a^T x \leq b,
              \end{align*} where $a \neq 0$.
        \item \begin{align*}
                  \minimize  & \quad  c^T x                                                      \\
                  \subjectto & \quad  \mathbf{1}^T x \leq \alpha                                 \\
                             & \quad  0 \cdot \mathbf{1} \preceq l \preceq x \preceq u. \newline
              \end{align*}
    \end{enumerate}

    \solution

    \textbf{Part (a):} Because of the half-space constraint, this problem is unbounded below, so the optimal solution is $-\infty$. \\

    \textbf{Part (b):} Splitting by components gives $c_i x_i$ subject to $0 \leq l_i \leq x_i \leq u_i$ for each $i = 1, 2, \ldots, n$.
    \[
        x_i^* =
        \begin{cases}
            l_i, & c_i > 0, \\
            u_i, & c_i < 0, \\
        \end{cases}
    \] where $\sum x_i^* \leq \alpha$.
\end{homeworkProblem}

\pagebreak

% ----- PROBLEM 8 ------ %
\begin{homeworkProblem}
    Find the dual optimization problems for the following optimization problems
    \begin{enumerate}[label=(\alph*), topsep=5pt, itemsep=5pt]
        \item \begin{align*}
                  \minimize  & \quad  c^T x                  \\
                  \subjectto & \quad  Ax = b                 \\
                             & \quad  l \preceq x \preceq u,
              \end{align*} where $x \in \R[n]$ is the variable, and $l, u \in \R[n]$
        \item Recall last time you wrote the following problem as an equivalent LP:
              \begin{align*}
                  \minimize  & \quad  \norm{x}_\infty              \\
                  \subjectto & \quad  \norm{Ax - b}_\infty \leq 1.
              \end{align*} where $A \in \R[m \times n]$ and $b \in \R[m]$. Find the dual problem for the equivalent LP.
        \item \begin{align*}
                  \minimize  & \quad \sum_{i=1}^n x_i \log x_i \\
                  \subjectto & \quad Ax \succeq b,             \\
                             & \quad c^T x = d,
              \end{align*} whose domain is $\R[n]_+$, i.e., $x_i > 0$ for all $i = 1, 2, \ldots, n$. Assume $c \succeq 0$.
        \item \begin{align*}
                  \minimize  & \quad  x^T Q x \\
                  \subjectto & \quad  Ax = b,
              \end{align*} where $Q$ is a positive definite matrix. \newline
    \end{enumerate}

    \solution

    \textbf{Part (a):}
    \begin{align*}
        g(\nu, \lambda_1, \lambda_2) & = \inf_{x \in \mathcal{D}} L (x, \nu, \lambda_1, \lambda_2)                                                                \\
                                     & = \inf_{x \in \mathcal{D}} \left( c^T x + \nu^T (Ax - b) + \lambda_1^T (x-u) + \lambda_2^T (l-x) \right)                   \\
                                     & = \inf_{x \in \mathcal{D}} \left(\lambda_2^T l - \lambda_1^T u - v^T b + (c^T + v^T A + \lambda_1^T - \lambda_2^T)x\right) \\
                                     & = \lambda_2^T l - \lambda_1^T u - v^T b
    \end{align*} where we maximize $g(\lambda, \nu_1, \nu_2)$ subject to $\lambda \succeq 0$.

    \pagebreak

    \textbf{Part (b):}
    \begin{align*}
        g(\lambda_1, \lambda_2, \lambda_3, \lambda_4) & = \inf_{t \in \mathcal{D}} L(t, \lambda_1, \lambda_2, \lambda_3, \lambda_4)                                                          \\
                                                      & = \inf_{t \in \mathcal{D}} \left( t + \lambda_1^T ( x- t) + \lambda_2^T (-t-x) + \lambda_3^T(Ax-b-t) + \lambda_4^T(-t -Ax+b) \right) \\
                                                      & = t(1 - \lambda_1^T - \lambda_2^T - \lambda_3^T - \lambda_4^T) + b^T(\lambda_3 - \lambda_4)
    \end{align*} where we maximize $g(\lambda_1, \lambda_2, \lambda_3, \lambda_4)$ subject to $\lambda \succeq 0$. \\

    \textbf{Part (c):}
    \begin{align*}
        g(\lambda, \nu) & = \inf_{x \in \mathcal{D}} L (x, \lambda, \nu)                                                             \\
                        & = \inf_{x \in \mathcal{D}} \left( \sum_{i=1}^n x_i \log x_i + \lambda^T (b - Ax) + \nu (c^T x - d) \right)
    \end{align*} where we maximize $g(\lambda, \nu)$. \\

    \textbf{Part (d):}
    \begin{align*}
        g(\nu) & = \inf_{x} L (x, \nu)                                                                     \\
               & = \inf_{x} \left( x^T Q x + \nu^T (Ax - b) \right)                                        \\
               & = (-\frac{1}{2} A^T \nu)^T Q (-\frac{1}{2} A^T \nu) + \nu^T (A(-\frac{1}{2} A^T \nu) - b)
    \end{align*} where we maximize $g(\nu)$.


\end{homeworkProblem}


\end{document}